% Seminar 3 - Page 3. Fragment for
% % Seminar 3 - Page 3. Fragment for
% % Seminar 3 - Page 3. Fragment for
% % Seminar 3 - Page 3. Fragment for
% \input{seminars/pages/seminar3-pagina3-en}; no preamble or forced page break.
% EDITORIAL NOTES:
% - The manuscript defines "countable" through a bijection from N; this
%   describes countably infinite sets. It then notes that a subset of such
%   a set is either finite or countable.
% - The exact symbol beside the subset remark is faint. Its mathematical
%   reading (finite or countable) is used below, provisionally.
% - Q ~ Z x Z is the example as written; no explicit bijection is supplied.
% - In the diagonal argument the manuscript uses a subscript scheme
%   y_{n1}, y_{n2}, ... and z_0,z_1,...; this is preserved.
% - At the contradiction line the letter is faint; z != y_n is inferred
%   from the immediately preceding construction.
% - The manuscript does not discuss the ambiguity of decimal expansions,
%   nor does it impose a convention for trailing nines. No additional
%   justification has been inserted into the transcription.
% - The product-of-injections note is separate from Exercise 1 in the source.
% - Exercise 2 is unfinished at the bottom of the manuscript page.

\subsection*{Seminar 3 - Page 3}

\noindent A set $A$ is called \emph{countable} if there exists a bijection
\[
  f:\mathbb N\longrightarrow A.
\]
\noindent\textbf{Examples.} $\mathbb N$, $\mathbb Z$,
$\mathbb Q\ (\sim\mathbb Z\times\mathbb Z)$.

\noindent If $A$ is countable and $B\subseteq A$, then $B$ is
finite or countable.

\medskip
\noindent\textbf{(1) $\mathbb R$ is not countable.}

\noindent\underline{Proof (by contradiction).}
If $\mathbb R$ were countable, then $[0,1)$ would be countable.
Suppose
\[
  [0,1)=\{y_n\mid n\in\mathbb N\},
\]
where the elements have decimal expansions
\[
\begin{aligned}
 y_0 &= 0,y_{01}y_{02}y_{03}\ldots,\\
 y_1 &= 0,y_{11}y_{12}y_{13}\ldots,\\
     &\ \vdots\\
 y_n &= 0,y_{n1}y_{n2}y_{n3}\ldots .
\end{aligned}
\]
Construct the number
\[
  z:=0,z_0z_1z_2\ldots z_t\ldots\in[0,1)
\]
by defining
\[
 z_0=\begin{cases}
       1+y_{01},&y_{01}\neq9,\\
       0,&y_{01}=9,
     \end{cases}
 \qquad
 z_t=\begin{cases}
       1+y_{t,t+1},&y_{t,t+1}\neq9,\\
       0,&y_{t,t+1}=9.
     \end{cases}
\]
Thus $z\neq y_n$ for every $n\in\mathbb N$, contradicting
$[0,1)=\{y_n\mid n\in\mathbb N\}$.
Consequently, $[0,1)$ is not countable and therefore
$\mathbb R$ is not countable.

\medskip
\noindent\emph{Note on products of countable sets.}
Let
\[
 A_1=\{x_1,x_2,\ldots,x_n,\ldots\},\qquad
 B_1=\{y_1,y_2,\ldots,y_n,\ldots\},
\]
with maps $f_1:A_1\to\mathbb N$ and $f_2:B_1\to\mathbb N$.
The manuscript records the product map
\[
\begin{aligned}
 f_1\times f_2 &: A_1\times B_1\longrightarrow
                    \mathbb N\times\mathbb N,\\
 (f_1\times f_2)(x_i,y_j)
   &=\bigl(f_1(x_i),f_2(y_j)\bigr).
\end{aligned}
\]

\medskip
\noindent\textbf{Exercise 2.} Show that there is no function
$f:\mathbb R\to\mathbb R$ such that
\[
 |f(x)-f(y)|>1\qquad(\forall x\neq y).
\]

\noindent\underline{Proof (by contradiction).}
Suppose that such a map $f:\mathbb R\to\mathbb R$ exists.
Then
\[
 x\neq y\ \Longrightarrow\ f(x)\neq f(y),
\]
so $f$ is injective. Let
\[
 \operatorname{im}(f)=\{f(x)\mid x\in\mathbb R\}.
\]
Since $f$ is injective,
$\operatorname{im}(f)\sim\mathbb R$, hence
$\operatorname{im}(f)$ is not countable.
On the other hand,
\[
 \operatorname{im}(f)
  =\operatorname{im}(f)\cap\mathbb R
  =\operatorname{im}(f)\cap
     \left(\bigcup_{n\in\mathbb Z}[n,n+1)\right)
  =\ \cdots
\]
% The argument continues on the next manuscript page.
; no preamble or forced page break.
% EDITORIAL NOTES:
% - The manuscript defines "countable" through a bijection from N; this
%   describes countably infinite sets. It then notes that a subset of such
%   a set is either finite or countable.
% - The exact symbol beside the subset remark is faint. Its mathematical
%   reading (finite or countable) is used below, provisionally.
% - Q ~ Z x Z is the example as written; no explicit bijection is supplied.
% - In the diagonal argument the manuscript uses a subscript scheme
%   y_{n1}, y_{n2}, ... and z_0,z_1,...; this is preserved.
% - At the contradiction line the letter is faint; z != y_n is inferred
%   from the immediately preceding construction.
% - The manuscript does not discuss the ambiguity of decimal expansions,
%   nor does it impose a convention for trailing nines. No additional
%   justification has been inserted into the transcription.
% - The product-of-injections note is separate from Exercise 1 in the source.
% - Exercise 2 is unfinished at the bottom of the manuscript page.

\subsection*{Seminar 3 - Page 3}

\noindent A set $A$ is called \emph{countable} if there exists a bijection
\[
  f:\mathbb N\longrightarrow A.
\]
\noindent\textbf{Examples.} $\mathbb N$, $\mathbb Z$,
$\mathbb Q\ (\sim\mathbb Z\times\mathbb Z)$.

\noindent If $A$ is countable and $B\subseteq A$, then $B$ is
finite or countable.

\medskip
\noindent\textbf{(1) $\mathbb R$ is not countable.}

\noindent\underline{Proof (by contradiction).}
If $\mathbb R$ were countable, then $[0,1)$ would be countable.
Suppose
\[
  [0,1)=\{y_n\mid n\in\mathbb N\},
\]
where the elements have decimal expansions
\[
\begin{aligned}
 y_0 &= 0,y_{01}y_{02}y_{03}\ldots,\\
 y_1 &= 0,y_{11}y_{12}y_{13}\ldots,\\
     &\ \vdots\\
 y_n &= 0,y_{n1}y_{n2}y_{n3}\ldots .
\end{aligned}
\]
Construct the number
\[
  z:=0,z_0z_1z_2\ldots z_t\ldots\in[0,1)
\]
by defining
\[
 z_0=\begin{cases}
       1+y_{01},&y_{01}\neq9,\\
       0,&y_{01}=9,
     \end{cases}
 \qquad
 z_t=\begin{cases}
       1+y_{t,t+1},&y_{t,t+1}\neq9,\\
       0,&y_{t,t+1}=9.
     \end{cases}
\]
Thus $z\neq y_n$ for every $n\in\mathbb N$, contradicting
$[0,1)=\{y_n\mid n\in\mathbb N\}$.
Consequently, $[0,1)$ is not countable and therefore
$\mathbb R$ is not countable.

\medskip
\noindent\emph{Note on products of countable sets.}
Let
\[
 A_1=\{x_1,x_2,\ldots,x_n,\ldots\},\qquad
 B_1=\{y_1,y_2,\ldots,y_n,\ldots\},
\]
with maps $f_1:A_1\to\mathbb N$ and $f_2:B_1\to\mathbb N$.
The manuscript records the product map
\[
\begin{aligned}
 f_1\times f_2 &: A_1\times B_1\longrightarrow
                    \mathbb N\times\mathbb N,\\
 (f_1\times f_2)(x_i,y_j)
   &=\bigl(f_1(x_i),f_2(y_j)\bigr).
\end{aligned}
\]

\medskip
\noindent\textbf{Exercise 2.} Show that there is no function
$f:\mathbb R\to\mathbb R$ such that
\[
 |f(x)-f(y)|>1\qquad(\forall x\neq y).
\]

\noindent\underline{Proof (by contradiction).}
Suppose that such a map $f:\mathbb R\to\mathbb R$ exists.
Then
\[
 x\neq y\ \Longrightarrow\ f(x)\neq f(y),
\]
so $f$ is injective. Let
\[
 \operatorname{im}(f)=\{f(x)\mid x\in\mathbb R\}.
\]
Since $f$ is injective,
$\operatorname{im}(f)\sim\mathbb R$, hence
$\operatorname{im}(f)$ is not countable.
On the other hand,
\[
 \operatorname{im}(f)
  =\operatorname{im}(f)\cap\mathbb R
  =\operatorname{im}(f)\cap
     \left(\bigcup_{n\in\mathbb Z}[n,n+1)\right)
  =\ \cdots
\]
% The argument continues on the next manuscript page.
; no preamble or forced page break.
% EDITORIAL NOTES:
% - The manuscript defines "countable" through a bijection from N; this
%   describes countably infinite sets. It then notes that a subset of such
%   a set is either finite or countable.
% - The exact symbol beside the subset remark is faint. Its mathematical
%   reading (finite or countable) is used below, provisionally.
% - Q ~ Z x Z is the example as written; no explicit bijection is supplied.
% - In the diagonal argument the manuscript uses a subscript scheme
%   y_{n1}, y_{n2}, ... and z_0,z_1,...; this is preserved.
% - At the contradiction line the letter is faint; z != y_n is inferred
%   from the immediately preceding construction.
% - The manuscript does not discuss the ambiguity of decimal expansions,
%   nor does it impose a convention for trailing nines. No additional
%   justification has been inserted into the transcription.
% - The product-of-injections note is separate from Exercise 1 in the source.
% - Exercise 2 is unfinished at the bottom of the manuscript page.

\subsection*{Seminar 3 - Page 3}

\noindent A set $A$ is called \emph{countable} if there exists a bijection
\[
  f:\mathbb N\longrightarrow A.
\]
\noindent\textbf{Examples.} $\mathbb N$, $\mathbb Z$,
$\mathbb Q\ (\sim\mathbb Z\times\mathbb Z)$.

\noindent If $A$ is countable and $B\subseteq A$, then $B$ is
finite or countable.

\medskip
\noindent\textbf{(1) $\mathbb R$ is not countable.}

\noindent\underline{Proof (by contradiction).}
If $\mathbb R$ were countable, then $[0,1)$ would be countable.
Suppose
\[
  [0,1)=\{y_n\mid n\in\mathbb N\},
\]
where the elements have decimal expansions
\[
\begin{aligned}
 y_0 &= 0,y_{01}y_{02}y_{03}\ldots,\\
 y_1 &= 0,y_{11}y_{12}y_{13}\ldots,\\
     &\ \vdots\\
 y_n &= 0,y_{n1}y_{n2}y_{n3}\ldots .
\end{aligned}
\]
Construct the number
\[
  z:=0,z_0z_1z_2\ldots z_t\ldots\in[0,1)
\]
by defining
\[
 z_0=\begin{cases}
       1+y_{01},&y_{01}\neq9,\\
       0,&y_{01}=9,
     \end{cases}
 \qquad
 z_t=\begin{cases}
       1+y_{t,t+1},&y_{t,t+1}\neq9,\\
       0,&y_{t,t+1}=9.
     \end{cases}
\]
Thus $z\neq y_n$ for every $n\in\mathbb N$, contradicting
$[0,1)=\{y_n\mid n\in\mathbb N\}$.
Consequently, $[0,1)$ is not countable and therefore
$\mathbb R$ is not countable.

\medskip
\noindent\emph{Note on products of countable sets.}
Let
\[
 A_1=\{x_1,x_2,\ldots,x_n,\ldots\},\qquad
 B_1=\{y_1,y_2,\ldots,y_n,\ldots\},
\]
with maps $f_1:A_1\to\mathbb N$ and $f_2:B_1\to\mathbb N$.
The manuscript records the product map
\[
\begin{aligned}
 f_1\times f_2 &: A_1\times B_1\longrightarrow
                    \mathbb N\times\mathbb N,\\
 (f_1\times f_2)(x_i,y_j)
   &=\bigl(f_1(x_i),f_2(y_j)\bigr).
\end{aligned}
\]

\medskip
\noindent\textbf{Exercise 2.} Show that there is no function
$f:\mathbb R\to\mathbb R$ such that
\[
 |f(x)-f(y)|>1\qquad(\forall x\neq y).
\]

\noindent\underline{Proof (by contradiction).}
Suppose that such a map $f:\mathbb R\to\mathbb R$ exists.
Then
\[
 x\neq y\ \Longrightarrow\ f(x)\neq f(y),
\]
so $f$ is injective. Let
\[
 \operatorname{im}(f)=\{f(x)\mid x\in\mathbb R\}.
\]
Since $f$ is injective,
$\operatorname{im}(f)\sim\mathbb R$, hence
$\operatorname{im}(f)$ is not countable.
On the other hand,
\[
 \operatorname{im}(f)
  =\operatorname{im}(f)\cap\mathbb R
  =\operatorname{im}(f)\cap
     \left(\bigcup_{n\in\mathbb Z}[n,n+1)\right)
  =\ \cdots
\]
% The argument continues on the next manuscript page.
; no preamble or forced page break.
% EDITORIAL NOTES:
% - The manuscript defines "countable" through a bijection from N; this
%   describes countably infinite sets. It then notes that a subset of such
%   a set is either finite or countable.
% - The exact symbol beside the subset remark is faint. Its mathematical
%   reading (finite or countable) is used below, provisionally.
% - Q ~ Z x Z is the example as written; no explicit bijection is supplied.
% - In the diagonal argument the manuscript uses a subscript scheme
%   y_{n1}, y_{n2}, ... and z_0,z_1,...; this is preserved.
% - At the contradiction line the letter is faint; z != y_n is inferred
%   from the immediately preceding construction.
% - The manuscript does not discuss the ambiguity of decimal expansions,
%   nor does it impose a convention for trailing nines. No additional
%   justification has been inserted into the transcription.
% - The product-of-injections note is separate from Exercise 1 in the source.
% - Exercise 2 is unfinished at the bottom of the manuscript page.

\subsection*{Seminar 3 - Page 3}

\noindent A set $A$ is called \emph{countable} if there exists a bijection
\[
  f:\mathbb N\longrightarrow A.
\]
\noindent\textbf{Examples.} $\mathbb N$, $\mathbb Z$,
$\mathbb Q\ (\sim\mathbb Z\times\mathbb Z)$.

\noindent If $A$ is countable and $B\subseteq A$, then $B$ is
finite or countable.

\medskip
\noindent\textbf{(1) $\mathbb R$ is not countable.}

\noindent\underline{Proof (by contradiction).}
If $\mathbb R$ were countable, then $[0,1)$ would be countable.
Suppose
\[
  [0,1)=\{y_n\mid n\in\mathbb N\},
\]
where the elements have decimal expansions
\[
\begin{aligned}
 y_0 &= 0,y_{01}y_{02}y_{03}\ldots,\\
 y_1 &= 0,y_{11}y_{12}y_{13}\ldots,\\
     &\ \vdots\\
 y_n &= 0,y_{n1}y_{n2}y_{n3}\ldots .
\end{aligned}
\]
Construct the number
\[
  z:=0,z_0z_1z_2\ldots z_t\ldots\in[0,1)
\]
by defining
\[
 z_0=\begin{cases}
       1+y_{01},&y_{01}\neq9,\\
       0,&y_{01}=9,
     \end{cases}
 \qquad
 z_t=\begin{cases}
       1+y_{t,t+1},&y_{t,t+1}\neq9,\\
       0,&y_{t,t+1}=9.
     \end{cases}
\]
Thus $z\neq y_n$ for every $n\in\mathbb N$, contradicting
$[0,1)=\{y_n\mid n\in\mathbb N\}$.
Consequently, $[0,1)$ is not countable and therefore
$\mathbb R$ is not countable.

\medskip
\noindent\emph{Note on products of countable sets.}
Let
\[
 A_1=\{x_1,x_2,\ldots,x_n,\ldots\},\qquad
 B_1=\{y_1,y_2,\ldots,y_n,\ldots\},
\]
with maps $f_1:A_1\to\mathbb N$ and $f_2:B_1\to\mathbb N$.
The manuscript records the product map
\[
\begin{aligned}
 f_1\times f_2 &: A_1\times B_1\longrightarrow
                    \mathbb N\times\mathbb N,\\
 (f_1\times f_2)(x_i,y_j)
   &=\bigl(f_1(x_i),f_2(y_j)\bigr).
\end{aligned}
\]

\medskip
\noindent\textbf{Exercise 2.} Show that there is no function
$f:\mathbb R\to\mathbb R$ such that
\[
 |f(x)-f(y)|>1\qquad(\forall x\neq y).
\]

\noindent\underline{Proof (by contradiction).}
Suppose that such a map $f:\mathbb R\to\mathbb R$ exists.
Then
\[
 x\neq y\ \Longrightarrow\ f(x)\neq f(y),
\]
so $f$ is injective. Let
\[
 \operatorname{im}(f)=\{f(x)\mid x\in\mathbb R\}.
\]
Since $f$ is injective,
$\operatorname{im}(f)\sim\mathbb R$, hence
$\operatorname{im}(f)$ is not countable.
On the other hand,
\[
 \operatorname{im}(f)
  =\operatorname{im}(f)\cap\mathbb R
  =\operatorname{im}(f)\cap
     \left(\bigcup_{n\in\mathbb Z}[n,n+1)\right)
  =\ \cdots
\]
% The argument continues on the next manuscript page.
